\documentclass[11pt, a4paper]{article}

% =====================
% PACKAGES
% =====================
\usepackage{graphicx}
\usepackage{amsmath}
\usepackage{amssymb}
\usepackage{geometry}
\usepackage{booktabs}
\usepackage{hyperref}
\usepackage{listings}
\usepackage{xcolor}
\usepackage{float}
\usepackage{caption}

% =====================
% CONFIG
% =====================
\geometry{margin=1in}
\definecolor{codegray}{rgb}{0.95,0.95,0.95}

\lstset{
    backgroundcolor=\color{codegray},
    basicstyle=\ttfamily\small,
    breaklines=true,
    frame=single,
    columns=fullflexible
}

% =====================
% TITLE
% =====================
\title{\textbf{Project RIFT-MNEMOS:}\\
\large Spatially Gated Memory Emergence in a Non-Equilibrium Antagonistic Substrate}

\author{
\textbf{Akshay John} \\
Independent Researcher \\
\texttt{Experiment Series: RIFT-003 / MNEMOS-V}
}

\date{January 2026}

\begin{document}
\maketitle

% =====================
% ABSTRACT
% =====================
\begin{abstract}
Modern artificial intelligence relies on optimization against explicit objectives.
In contrast, biological intelligence emerges from systems forced to persist under continuous internal and environmental stress.
We present \textbf{RIFT-MNEMOS}, a non-optimizing, antagonistically regulated substrate in which memory arises as irreversible topological deformation.
We demonstrate that naive scarring mechanisms fail due to global callus formation and catastrophic interference.
By introducing spatial novelty gating, the system selectively consolidates episodic trauma while remaining robust to background homeostatic noise.
This work establishes a reproducible physical mechanism for distributed memory formation without loss functions, backpropagation, or symbolic supervision.
\end{abstract}

% =====================
% INTRODUCTION
% =====================
\section{Introduction}

Transformers and related architectures are fundamentally reactive systems: computation begins only upon prompt invocation.
Biological systems, in contrast, exist continuously and must actively resist entropy.
Project RIFT explores whether intelligence prerequisites can emerge from \textit{forced persistence}, rather than task-driven optimization.

This paper does not claim intelligence.
It documents the construction and failure-driven refinement of a memory-capable substrate.

% =====================
% ARCHITECTURE
% =====================
\section{System Architecture}

RIFT consists of three antagonistic components:

\begin{enumerate}
    \item \textbf{Tectonic Substrate}: A $48 \times 48 \times 32$ continuous grid updated via local stress diffusion.
    \item \textbf{Fracture Law}: Local fracture occurs when stress exceeds elasticity, dissipating energy while hardening structure.
    \item \textbf{Antagonist Ego}: A global regulator enforcing non-equilibrium:
    \begin{itemize}
        \item Low fracture $\rightarrow$ inject noise (PRESSURIZE)
        \item High fracture $\rightarrow$ damp state (CONSTRAIN)
    \end{itemize}
\end{enumerate}

The system is always active. There is no inference mode.

% =====================
% FAILURE MODES
% =====================
\section{Observed Failure Modes}

\subsection{Thermodynamic Explosion}
Unbounded stress amplification resulted in numerical divergence.

\subsection{Thermodynamic Coma}
Excessive damping collapsed the system into inert equilibrium.

\subsection{Global Callus Formation}
Single-layer elasticity caused background Ego noise to harden the entire substrate, blocking future learning.

\subsection{Catastrophic Interference}
Sequential trauma erased prior memory traces without spatial isolation.

% =====================
% EPISODIC MEMORY
% =====================
\section{Experiment I: Episodic Trauma Adaptation}

\subsection{Protocol}
A localized impulse ($F=6.0$) was injected every 50 timesteps.

\subsection{Result}
Peak fracture at impact decreased monotonically across exposures, demonstrating structural adaptation.

\begin{lstlisting}
Impact 1: 0.0644
Impact 6: 0.0092
Impact 12: 0.0036
\end{lstlisting}

% =====================
% CAPACITY TESTS
% =====================
\section{Experiment II: Spatial Capacity (MNEMOS-II/III)}

Sequential trauma at 9 sites revealed severe interference.
Maximum retention without intervention was 4/9 sites.

% =====================
% DUAL TIMESCALE FAILURE
% =====================
\section{Experiment III: Dual-Timescale Failure}

Separating fast and slow elasticity alone caused \textit{homeostatic insulation}, preventing trauma from reaching long-term memory.

% =====================
% SPATIAL GATING
% =====================
\section{Experiment IV: Spatial Gating (MNEMOS-V)}

\subsection{Method}
Memory consolidation was gated by spatial sharpness:

\[
G(x,y) = |\sigma(x,y) - \mu_{\text{local}}|
\]

Slow memory updated only if $G(x,y) > \theta$.

\subsection{Result}

\begin{lstlisting}
Sites Retained: 7 / 9
Blocked Sites: Stress-diffusion overlap
\end{lstlisting}

\subsection{Interpretation}
Memory is physically coupled, not discretely addressable.

% =====================
% DISCUSSION
% =====================
\section{Discussion}

RIFT-MNEMOS demonstrates:
\begin{itemize}
    \item Memory without optimization
    \item Structural persistence under antagonism
    \item Resolution of stability-plasticity via geometry
\end{itemize}

It does not demonstrate reasoning, abstraction, or intelligence.

% =====================
% FUTURE WORK
% =====================
\section{Future Directions}

Future research will explore:
\begin{itemize}
    \item Predictive dynamics based on scar geometry
    \item Non-symbolic input modalities
    \item Emergent internal world models
\end{itemize}

\section{Conclusion}

We have constructed a physical memory substrate.
Intelligence remains an open question.

\end{document}
